% GENERATED FILE. Edit canonical lessons, then run npm run generate:books.
% canonical-source-hash: fnv1a64:da4de4deb3df08e1
% canonical-lessons: JA-C267-maigo, JA-C267-yobidasu, JA-C267-housou, JA-C267-nakigoe, JA-C267-hotto, JA-R267-review-262-267

\chapter{A Lost Child}
\label{ch:ja-a-lost-child}
\hlchaptermodality{267}

\begin{chapteropening}
\textbf{By the end of this chapter:} \emph{I can say who went missing and how they were found.}

Complete the last lesson of chapter 267: I can say who went missing and how they were found.
\end{chapteropening}

\section[\texorpdfstring{\ja{まいご} (maigo)}{まいご (maigo)}]{\ja{まいご} (maigo) — a lost child}

\label{lesson:JA-C267-maigo}



\begin{quote}
Before the new one: say the Japanese for at last, then the Japanese for unharmed.
\end{quote}

\subsection*{You'll want to know: \ja{まいご}}
\textbf{\ja{まいご}} — \emph{maigo} — \textquotedblleft{}a lost child\textquotedblright{}.

\begin{quote}
\textbf{\ja{もし} \ja{まいごに} \ja{なったら、こうばんに} \ja{いきます。}} — \emph{moshi maigo ni nattara, kōban ni ikimasu} — \textquotedblleft{}If I get lost, I go to the police box.\textquotedblright{}
\end{quote}

\begin{cousinweb}
\textbf{\ja{まいご}} (\emph{maigo}) is a native Japanese word (wago), inherited rather than borrowed from Chinese, built from \emph{mai}, \textquotedblleft{}losing the way\textquotedblright{} (from \emph{mayou}), and \emph{go}, \textquotedblleft{}child\textquotedblright{} (from \emph{ko}) (JMdict, KANJIDIC2).
\end{cousinweb}

\begin{grammarlens}[title={What you've built}]
One more word for telling what happened.
\end{grammarlens}

\subsection*{Your turn}
\begin{itemize}
\raggedright
  \item \emph{Say it:} \emph{maigo}
  \item \emph{Say it:} \emph{maigo}, once more
  \item \emph{Say it:} \emph{moshi maigo ni nattara, kōban ni ikimasu} — If I get lost, I go to the police box.
  \item \emph{Recall:} say the Japanese for at last, then the Japanese for unharmed, then say \emph{maigo} again
\end{itemize}

\begin{tcolorbox}[breakable,colback=teal!4,colframe=teal!35!black,title={Before you move on}]
What does \ja{まいご} mean? (\textquotedblleft{}A lost child\textquotedblright{}.) Say it once more.
\end{tcolorbox}
\section[\texorpdfstring{\ja{よびだす} (yobidasu)}{よびだす (yobidasu)}]{\ja{よびだす} (yobidasu) — to call (someone) out; to page}

\label{lesson:JA-C267-yobidasu}



\begin{quote}
Before the new one: say the Japanese for unharmed, then the Japanese for a lost child.
\end{quote}

\subsection*{You'll want to know: \ja{よびだす}}
\textbf{\ja{よびだす}} — \emph{yobidasu} — \textquotedblleft{}to call (someone) out; to page\textquotedblright{}.

\begin{quote}
\textbf{\ja{えきいんが} \ja{ははを} \ja{よびだしました。}} — \emph{ekiin ga haha o yobidashimashita} — \textquotedblleft{}A station attendant paged my mother.\textquotedblright{}
\end{quote}

Polite \textbf{\ja{よびだします}}, \emph{yobidashimasu}; a godan verb. Past (polite): \textbf{\ja{せんせいが} \ja{ははを} \ja{よびだしました。}}, \emph{sensei ga haha o yobidashimashita}.

\begin{cousinweb}
\textbf{\ja{よびだす}} (\emph{yobidasu}) is a native Japanese word (wago), inherited rather than borrowed from Chinese, built from \emph{yobi}, \textquotedblleft{}calling\textquotedblright{} (from \emph{yobu}), and \emph{dasu}, \textquotedblleft{}to put out\textquotedblright{} (JMdict, KANJIDIC2).
\end{cousinweb}

\begin{grammarlens}[title={What you've built}]
One more word for telling what happened.
\end{grammarlens}

\subsection*{Your turn}
\begin{itemize}
\raggedright
  \item \emph{Say it:} \emph{yobidasu}
  \item \emph{Say it:} \emph{yobidasu}, once more
  \item \emph{Say it:} \emph{ekiin ga haha o yobidashimashita} — A station attendant paged my mother.
  \item \emph{Recall:} say the Japanese for unharmed, then the Japanese for a lost child, then say \emph{yobidasu} again
\end{itemize}

\begin{tcolorbox}[breakable,colback=teal!4,colframe=teal!35!black,title={Before you move on}]
What does \ja{よびだす} mean? (\textquotedblleft{}To call (someone) out; to page\textquotedblright{}.) Say it once more.
\end{tcolorbox}
\section[\texorpdfstring{\ja{ほうそう} (hōsō)}{ほうそう (hōsō)}]{\ja{ほうそう} (hōsō) — an announcement; a broadcast}

\label{lesson:JA-C267-housou}



\begin{quote}
Before the new one: say the Japanese for a lost child, then the Japanese for to call (someone) out.
\end{quote}

\subsection*{You'll want to know: \ja{ほうそう}}
\textbf{\ja{ほうそう}} — \emph{hōsō} — \textquotedblleft{}an announcement; a broadcast\textquotedblright{}.

\begin{quote}
\textbf{\ja{えきで} \ja{まいごの} \ja{ほうそうが} \ja{ありました。}} — \emph{eki de maigo no hōsō ga arimashita} — \textquotedblleft{}At the station they announced a lost child.\textquotedblright{}
\end{quote}

\begin{cousinweb}
\textbf{\ja{ほうそう}} (\emph{hōsō}) is Sino-Japanese (kango): two Chinese characters read with their on'yomi, readings borrowed from Chinese. They are \emph{hō}, \textquotedblleft{}release\textquotedblright{} (Middle Chinese \emph{pjangH}) and \emph{sō}, \textquotedblleft{}send\textquotedblright{} (\emph{suwngH}) (KANJIDIC2; Middle Chinese in Baxter's transcription).
\end{cousinweb}

\begin{grammarlens}[title={What you've built}]
One more word for telling what happened.
\end{grammarlens}

\subsection*{Your turn}
\begin{itemize}
\raggedright
  \item \emph{Say it:} \emph{hōsō}
  \item \emph{Say it:} \emph{hōsō}, once more
  \item \emph{Say it:} \emph{eki de maigo no hōsō ga arimashita} — At the station they announced a lost child.
  \item \emph{Recall:} say the Japanese for a lost child, then the Japanese for to call (someone) out, then say \emph{hōsō} again
\end{itemize}

\begin{tcolorbox}[breakable,colback=teal!4,colframe=teal!35!black,title={Before you move on}]
What does \ja{ほうそう} mean? (\textquotedblleft{}An announcement; a broadcast\textquotedblright{}.) Say it once more.
\end{tcolorbox}
\section[\texorpdfstring{\ja{なきごえ} (nakigoe)}{なきごえ (nakigoe)}]{\ja{なきごえ} (nakigoe) — a cry; the sound of crying}

\label{lesson:JA-C267-nakigoe}



\begin{quote}
Before the new one: say the Japanese for to call (someone) out, then the Japanese for an announcement.
\end{quote}

\subsection*{You'll want to know: \ja{なきごえ}}
\textbf{\ja{なきごえ}} — \emph{nakigoe} — \textquotedblleft{}a cry; the sound of crying\textquotedblright{}.

\begin{quote}
\textbf{\ja{こどもの} \ja{なきごえが} \ja{しました。}} — \emph{kodomo no nakigoe ga shimashita} — \textquotedblleft{}We heard a child crying.\textquotedblright{}
\end{quote}

\begin{cousinweb}
\textbf{\ja{なきごえ}} (\emph{nakigoe}) is a native Japanese word (wago), inherited rather than borrowed from Chinese, built from \emph{naki}, \textquotedblleft{}crying\textquotedblright{} (from \emph{naku}), and \emph{goe}, \textquotedblleft{}voice\textquotedblright{} (from \emph{koe}) (JMdict, KANJIDIC2).
\end{cousinweb}

\begin{grammarlens}[title={What you've built}]
One more word for telling what happened.
\end{grammarlens}

\subsection*{Your turn}
\begin{itemize}
\raggedright
  \item \emph{Say it:} \emph{nakigoe}
  \item \emph{Say it:} \emph{nakigoe}, once more
  \item \emph{Say it:} \emph{kodomo no nakigoe ga shimashita} — We heard a child crying.
  \item \emph{Recall:} say the Japanese for to call (someone) out, then the Japanese for an announcement, then say \emph{nakigoe} again
\end{itemize}

\begin{tcolorbox}[breakable,colback=teal!4,colframe=teal!35!black,title={Before you move on}]
What does \ja{なきごえ} mean? (\textquotedblleft{}A cry; the sound of crying\textquotedblright{}.) Say it once more.
\end{tcolorbox}
\section[\texorpdfstring{\ja{ほっと} (hotto)}{ほっと (hotto)}]{\ja{ほっと} (hotto) — with relief}

\label{lesson:JA-C267-hotto}



\begin{quote}
Before the new one: say the Japanese for an announcement, then the Japanese for a cry.
\end{quote}

\subsection*{You'll want to know: \ja{ほっと}}
\textbf{\ja{ほっと}} — \emph{hotto} — \textquotedblleft{}with relief\textquotedblright{}.

\begin{quote}
\textbf{\ja{こどもが} \ja{みつかって、ほっと} \ja{しました。}} — \emph{kodomo ga mitsukatte, hotto shimashita} — \textquotedblleft{}The child was found, and we breathed again.\textquotedblright{}
\end{quote}

\begin{cousinweb}
\textbf{\ja{ほっと}} (\emph{hotto}) is a native Japanese mimetic word, inherited rather than borrowed from Chinese, and its older origin is not settled (JMdict).
\end{cousinweb}

\begin{grammarlens}[title={What you've built}]
One more word for telling what happened.
\end{grammarlens}

\subsection*{Your turn}
\begin{itemize}
\raggedright
  \item \emph{Say it:} \emph{hotto}
  \item \emph{Say it:} \emph{hotto}, once more
  \item \emph{Say it:} \emph{kodomo ga mitsukatte, hotto shimashita} — The child was found, and we breathed again.
  \item \emph{Recall:} say the Japanese for an announcement, then the Japanese for a cry, then say \emph{hotto} again
\end{itemize}

\begin{tcolorbox}[breakable,colback=teal!4,colframe=teal!35!black,title={Before you move on}]
What does \ja{ほっと} mean? (\textquotedblleft{}With relief\textquotedblright{}.) Say it once more.
\end{tcolorbox}
\section[(dialogue)]{Looking back: from An Accident by the Station to A Lost Child}

\label{lesson:JA-R267-review-262-267}



\begin{quote}
Say the Japanese for an accident and to be saved.
\end{quote}

\subsection*{Your turn}
Say these aloud:

\begin{itemize}
\raggedright
  \item an accident, to be saved, to run away, suddenly, damage
  \item an ambulance, to call, to shout, to panic, at once
  \item to fall, to bump (something against), a bandage, very, for a while
  \item smoke, to spread, a fire engine, an escape route, to rush to the scene
  \item to drown, to jump in, to tug, at last, unharmed
  \item a lost child, to call (someone) out, an announcement, a cry, with relief
\end{itemize}

\begin{tcolorbox}[breakable,colback=teal!4,colframe=teal!35!black,title={Before you move on}]
An accident? (\textbf{\ja{じこ}}, \emph{jiko}.) To be saved? (\textbf{\ja{たすかる}}, \emph{tasukaru}.) To run away? (\textbf{\ja{にげる}}, \emph{nigeru}.) Suddenly? (\textbf{\ja{とつぜん}}, \emph{totsuzen}.) Damage? (\textbf{\ja{ひがい}}, \emph{higai}.) An ambulance? (\textbf{\ja{きゅうきゅうしゃ}}, \emph{kyūkyūsha}.) To call? (\textbf{\ja{よぶ}}, \emph{yobu}.) To shout? (\textbf{\ja{さけぶ}}, \emph{sakebu}.) To panic? (\textbf{\ja{あわてる}}, \emph{awateru}.) At once? (\textbf{\ja{すぐ}}, \emph{sugu}.) To fall? (\textbf{\ja{おちる}}, \emph{ochiru}.) To bump (something against)? (\textbf{\ja{ぶつける}}, \emph{butsukeru}.) A bandage? (\textbf{\ja{ほうたい}}, \emph{hōtai}.) Very? (\textbf{\ja{とても}}, \emph{totemo}.) For a while? (\textbf{\ja{しばらく}}, \emph{shibaraku}.) Smoke? (\textbf{\ja{けむり}}, \emph{kemuri}.) To spread? (\textbf{\ja{ひろがる}}, \emph{hirogaru}.) A fire engine? (\textbf{\ja{しょうぼうしゃ}}, \emph{shōbōsha}.) An escape route? (\textbf{\ja{にげみち}}, \emph{nigemichi}.) To rush to the scene? (\textbf{\ja{かけつける}}, \emph{kaketsukeru}.) To drown? (\textbf{\ja{おぼれる}}, \emph{oboreru}.) To jump in? (\textbf{\ja{とびこむ}}, \emph{tobikomu}.) To tug? (\textbf{\ja{ひっぱる}}, \emph{hipparu}.) At last? (\textbf{\ja{やっと}}, \emph{yatto}.) Unharmed? (\textbf{\ja{ぶじ}}, \emph{buji}.) A lost child? (\textbf{\ja{まいご}}, \emph{maigo}.) To call (someone) out? (\textbf{\ja{よびだす}}, \emph{yobidasu}.) An announcement? (\textbf{\ja{ほうそう}}, \emph{hōsō}.) A cry? (\textbf{\ja{なきごえ}}, \emph{nakigoe}.) With relief? (\textbf{\ja{ほっと}}, \emph{hotto}.)
\end{tcolorbox}
